% ============================================================================
% Fichier  : partie2/P02_C06_conduire_audit_cyber.tex
% Titre    : Conduire un Audit de Cybersécurité
% Auteur   : MINKA MI NGUIDJOÏ Thierry Emmanuel
% Version  : 1.0 -- novembre 2026
% Corrigés : corriges/Corr_C06.tex
% ============================================================================
% Sources : ISO 19011:2018 (sept principes, programme et conduite d'audit :
% paraphrasés) ; ISO/IEC 27007:2020 ; IIA, normes mondiales d'audit interne
% (2024) ; ITAF 4e éd. (2020) ; NIST SP 800-53A rév. 5 (méthodes examiner,
% interroger, tester) ; loi n° 2010/012, art. 4 point 9, art. 13 al. 1,
% art. 32 al. 1 et 2 (texte français officiel, Ann. D d'ACSI).
% Faits du cas (bible) : BANQUE SANAGA, Mme EKAMBI (audit interne), plan
% A1 à A8, état de septembre 2027 (5/2/3), mission sur les agences
% programmée pour 2028 (Ex. 1.4), GAB de 14 agences sous un système sans
% correctifs, droits d'administration locale des chefs d'agence dans 31
% agences (C03).
% Faits nouveaux (reportés dans la bible) : mission « sécurité du réseau
% d'agences », février à avril 2028 ; échantillon de 8 agences sur 38 ;
% 1 240 comptes ; 8 tests T1 à T8 : 3 OUI / 3 NON / 2 NEK ; Kribi échoue à
% l'exercice de mode dégradé ; 31 agences sur 38 journalisées ; 3 agences
% sur 8 gardent l'administration locale ; le prestataire des GAB refuse
% l'accès ; contestation de la DSI à la réunion de clôture.
% ============================================================================

\chapter{Conduire un Audit de Cybersécurité}
\label{chap:audit}
\chaptermark{Conduire un audit}

\epigraph{%
    « L'auditeur ne croit pas.\\
    Il vérifie, ou il dit qu'il n'a pas pu. »%
}{Principe du cours}

\niveauChapitre{Fondamental}{%
    Maîtrise du Chapitre~\ref{chap:grc}%
}

\begin{boxobjectifs}
\begin{itemize}
    \item Définir l'audit, ses critères, ses preuves, ses constats et ses
          conclusions.
    \item Appliquer les sept principes de l'ISO~19011 et situer
          l'ISO/IEC~27007.
    \item Distinguer audit interne, audit de certification, audit
          contractuel et audit réglementaire.
    \item Décrire le cycle d'une mission, du mandat au suivi, et organiser
          le contradictoire.
    \item Choisir une technique de collecte et apprécier la fiabilité d'une
          preuve.
    \item Échantillonner, et dire ce qu'un échantillon permet de conclure.
    \item Coter un test \OUI{}, \NON{} ou \NEK{}, et traiter une limitation
          de l'étendue.
\end{itemize}
\end{boxobjectifs}

\bigskip

% ============================================================================
\section*{Avant de commencer : « Les actions sont faites »}
% ============================================================================

Douala, février 2028. Le plan de traitement de R-AG-01 est déclaré achevé.
Le rapport de la RSSI au comité des risques l'affirme~: second facteur,
cloisonnement, sauvegardes, mode dégradé, sensibilisation, journalisation,
tout est en place\footnote{Comme dans tout le cours, la BANQUE SANAGA, ses
dirigeants, ses prestataires, les dates et les montants sont fictifs~; les
normes et les textes cités sont exacts.}. Le comité ne doute pas de la
bonne foi de Mme~ABENA. Mais il a appris, en 2026, ce que vaut une
assurance qu'on ne vérifie pas. Il confirme donc la mission inscrite au
programme de l'audit interne~: la sécurité du réseau d'agences.

Mme~EKAMBI, directrice de l'audit interne, sait que la mission sera
regardée de près. Si elle conclut que tout va bien, le comité réduira la
surveillance~; si elle se trompe, la prochaine attaque le dira. Elle doit
donc pouvoir montrer, pour chaque conclusion, la preuve qui la fonde, ou
dire clairement qu'elle ne l'a pas obtenue.

\medskip

Ce chapitre décrit la conduite d'un audit de cybersécurité, de l'audit
interne à l'audit réglementaire. Le Chapitre~\ref{chap:conclure} apprendra
ensuite à décomposer l'objet audité et à conclure.

\begin{boxalert}
\textbf{Violation typique.} Conclure qu'une mesure est en place sur la
foi d'un rapport de l'audité, d'une capture d'écran qu'il a fournie, ou
d'un entretien.

\smallskip
\textbf{Règle~:} une déclaration n'est pas une preuve de fonctionnement.
L'auditeur obtient lui-même la preuve, ou une preuve produite par le
système sous ses yeux~; à défaut, il cote \NEK{} et le dit.
\end{boxalert}

\separateur

% ============================================================================
\section{Ce qu'est un Audit}
\label{sec:c06-definir}
% ============================================================================

\subsection{Cinq notions}

\begin{boxdef}
\termecle{Audit}\index{audit}~: processus méthodique, indépendant et
documenté, qui consiste à obtenir des preuves objectives et à les évaluer
de manière objective pour déterminer dans quelle mesure des critères sont
satisfaits~\cite{iso19011_2018}.

\medskip
\termecle{Critère d'audit}\index{critère d'audit}~: ensemble d'exigences
qui sert de référence~: une norme, une politique, une loi, un contrat, un
plan de traitement.

\medskip
\termecle{Preuve d'audit}\index{preuve d'audit}~: enregistrement,
énoncé de fait ou autre information, vérifiable et pertinent au regard des
critères.

\medskip
\termecle{Constat d'audit}\index{constat d'audit}~: résultat de
l'évaluation des preuves au regard des critères~: conformité,
non-conformité, ou impossibilité de conclure.

\medskip
\termecle{Conclusion d'audit}\index{conclusion d'audit}~: résultat de
l'audit, tiré de l'ensemble des constats, au regard des objectifs de la
mission.
\end{boxdef}
{\small\itshape Définitions reformulées d'après l'ISO~19011:2018~; la
troisième issue du constat est propre au cours.}

\medskip

L'ISO~19011 envisage deux issues pour un constat~: conforme ou non
conforme. Le cours en ajoute une troisième, \NEK{}~: la preuve n'a pas pu
être obtenue, ou ne permet pas de trancher. Ce n'est pas une
complaisance~; c'est une information, souvent la plus importante du
rapport (Chapitre~\ref{chap:conclure}).

\subsection{Sept principes}

L'ISO~19011:2018 pose sept principes, qui s'appliquent à tout audit de
système de management~\cite{iso19011_2018}.

\begin{tableaucours}{L{3.6cm}L{8.6cm}}{%
    \tableauHeader{Principe} &
    \tableauHeader{Ce qu'il signifie pour l'auditeur en cybersécurité}
}
Intégrité & Honnêteté, diligence, compétence~; refus de toute pression sur
les constats \\
Présentation impartiale & Rapporter fidèlement, y compris les obstacles
rencontrés et les divergences non résolues \\
Conscience professionnelle & Le soin qu'exige l'importance de la mission
et la confiance des commanditaires \\
Confidentialité & Protéger les informations obtenues, en particulier les
faiblesses découvertes \\
Indépendance & Ne pas auditer son propre travail~; rester objectif \\
Approche fondée sur la preuve & Conclure sur des preuves vérifiables,
obtenues sur un échantillon raisonné \\
Approche fondée sur le risque & Concentrer l'effort là où le risque, pour
la mission et pour l'audité, est le plus grand \\
\end{tableaucours}
\vspace{-0.8em}
\begin{center}{\small\itshape Tableau~6.1. Les sept principes de
l'audit (d'après l'ISO~19011:2018)}\end{center}

\begin{boxinfo}
\textbf{Les autres référentiels.} L'ISO/IEC~27007:2020 complète
l'ISO~19011 pour l'audit des systèmes de management de la sécurité de
l'information~\cite{iso27007_2020}~; le Chapitre~\ref{chap:smsi} l'emploie.
Les auditeurs internes appliquent en outre les normes mondiales d'audit
interne de l'IIA, publiées en 2024~\cite{iia2024gias}, et les auditeurs des
systèmes d'information le cadre ITAF de l'ISACA~\cite{isaca2020itaf}.
\end{boxinfo}

\subsection{Quatre familles d'audit}

\begin{tableaucours}{L{2.6cm}L{3.4cm}L{3.0cm}L{3.2cm}}{%
    \tableauHeader{Famille} &
    \tableauHeader{Commanditaire} &
    \tableauHeader{Critères} &
    \tableauHeader{Exemple du cours}
}
Interne (première partie) & L'organisation elle-même, par son audit
interne & Ses politiques, ses plans, les normes qu'elle applique &
Mission de Mme~EKAMBI sur les agences \\
Contractuel (seconde partie) & Un client ou un partenaire & Le contrat &
PHARMACAM chez LiZbiZ \\
Certification (tierce partie) & L'organisation, auprès d'un organisme
accrédité & La norme certifiable & CERTIQUAL AFRIQUE, ISO/IEC~27001 \\
Réglementaire & La loi, par une autorité & Les textes & L'audit de
sécurité obligatoire de l'ANPP \\
\end{tableaucours}
\vspace{-0.8em}
\begin{center}{\small\itshape Tableau~6.2. Les familles d'audit}\end{center}

Un test d'intrusion n'est pas un audit au sens strict~: il cherche des
chemins d'attaque, sans évaluer l'ensemble des critères. Il peut fournir
des preuves à un audit, comme le test A3 du plan de la banque. Il est
enseigné dans l'UE MSN-2S3-SEC03.

% ============================================================================
\section{Le Cycle de la Mission}
\label{sec:c06-cycle}
% ============================================================================

\subsection{Programme et mission}

L'ISO~19011 distingue le \emph{programme d'audit}, qui organise les audits
sur plusieurs années selon les risques, et la \emph{mission}, qui en
réalise un. Le programme de l'audit interne de la BANQUE SANAGA n'avait
pas couvert la sécurité en 2024 et 2025 (Chapitre~\ref{chap:grc},
Tableau~1.7)~; celui de 2028 l'inscrit en tête, parce que le risque
résiduel de R-AG-01 repose sur des mesures nouvelles.

\begin{figure}[H]
\centering
\begin{tikzpicture}[
    etape/.style={draw=CouleurPrimaire, thick, rounded corners,
                  fill=CouleurPrimaire!8, text width=2.0cm,
                  minimum height=1.2cm, align=center, font=\scriptsize},
    fleche/.style={->, thick, CouleurSecondaire}]
    \node[etape] (M) at (0,0) {\textbf{Mandat}\\objectifs, périmètre, critères};
    \node[etape] (P) at (2.5,0) {\textbf{Préparation}\\analyse de risque, programme de travail};
    \node[etape] (T) at (5.0,0) {\textbf{Terrain}\\collecte et évaluation des preuves};
    \node[etape] (C) at (7.5,0) {\textbf{Contradictoire}\\réunion de clôture, réponses};
    \node[etape] (R) at (10.0,0) {\textbf{Rapport}\\constats, conclusions};
    \node[etape] (S) at (5.0,-2.0) {\textbf{Suivi}\\vérification des actions};
    \draw[fleche] (M) -- (P); \draw[fleche] (P) -- (T);
    \draw[fleche] (T) -- (C); \draw[fleche] (C) -- (R);
    \draw[fleche] (R.south) |- (S.east);
    \draw[fleche, dashed] (S.west) -| (M.south);
\end{tikzpicture}
\caption{Le cycle d'une mission d'audit}
\label{fig:c06-cycle}
\end{figure}

\subsection{Le mandat}

\begin{boxscenario}{Opération SANAGA. La lettre de mission}
\textbf{Commanditaire~:} le comité des risques du conseil, auquel l'audit
interne rend compte.

\medskip
\textbf{Objectif~:} établir si les mesures du plan de traitement de
R-AG-01, et les mesures sur les agences issues de l'appréciation de
décembre 2026, sont en place et fonctionnent.

\medskip
\textbf{Périmètre~:} les 38 agences, les comptes et les postes des agents,
les espaces de fichiers, les sauvegardes des serveurs d'agence, les
automates de retrait (GAB)~; période~: de septembre 2027 à mars 2028.

\medskip
\textbf{Critères~:} la politique de sécurité de la banque~; le plan de
traitement (actions A1 à A8, preuves attendues)~; les mesures
correspondantes de l'ISO/IEC~27002:2022.

\medskip
\textbf{Équipe~:} Mme~EKAMBI et deux auditeurs internes, avec l'appui
d'un auditeur technique externe pour les extractions. Aucun membre de
l'équipe n'a participé à la mise en œuvre du plan.

\medskip
\textbf{Accès~:} la lettre de mission, signée par le président du comité,
donne à l'équipe l'accès à tous les systèmes, documents et personnes
nécessaires, y compris chez les prestataires dans la limite des contrats.
\end{boxscenario}

La dernière phrase compte. Sans elle, l'auditeur dépend de la bonne volonté
de chacun. Avec elle, un refus d'accès devient un constat, que le comité
doit traiter.

\subsection{La préparation}

L'équipe analyse les risques de la mission~: où une erreur de conclusion
serait la plus grave, où la preuve sera la plus difficile à obtenir. Elle
en tire un \emph{programme de travail}~: pour chaque critère, le test, la
technique, l'échantillon, la preuve attendue.

\begin{tableaucours}{cL{3.4cm}L{4.4cm}L{3.4cm}}{%
    \tableauHeader{Réf.} &
    \tableauHeader{Critère} &
    \tableauHeader{Test} &
    \tableauHeader{Preuve attendue}
}
T1 & Second facteur sur tous les comptes (A1) & Extraction exhaustive des
1\,240 comptes~; essai de connexion sans second facteur & Extraction,
procès-verbal d'essai \\
T2 & Cloisonnement par agence (A2) & Essais d'accès d'une agence vers les
espaces des autres & Procès-verbal d'essai \\
T3 & Sauvegardes hors ligne testées (A4) & Restauration observée d'un
serveur choisi par l'auditeur & Procès-verbal, empreinte des fichiers \\
T4 & Mode dégradé exercé (A5) & Exercice inopiné dans deux agences &
Compte rendu, chronométrage \\
T5 & Sensibilisation des agents (A6) & Rapprochement des listes de présence
et des effectifs & Listes, fichier du personnel \\
T6 & Journalisation centralisée (A8) & Extraction du collecteur~: agences
qui envoient des journaux & Extraction datée \\
T7 & Fin de l'administration locale des chefs d'agence & Extraction des
groupes locaux des postes & Extraction \\
T8 & Correctifs des GAB & Revue des configurations auprès du prestataire &
Configurations \\
\end{tableaucours}
\vspace{-0.8em}
\begin{center}{\small\itshape Tableau~6.3. Le programme de travail de la
mission}\end{center}

% ============================================================================
\section{Collecter et Apprécier la Preuve}
\label{sec:c06-preuve}
% ============================================================================

\subsection{Trois techniques, et leurs variantes}

Le NIST décrit trois méthodes d'évaluation des mesures de sécurité~:
\emph{examiner} (des documents, des configurations, des journaux),
\emph{interroger} (des personnes), \emph{tester} (faire fonctionner la
mesure dans des conditions définies)~\cite{nist2022sp80053a}. Les auditeurs
y ajoutent l'\emph{observation} (regarder une opération se faire) et la
\emph{ré-exécution} (refaire soi-même un contrôle pour voir s'il donne le
même résultat).

\subsection{La fiabilité de la preuve}

Toutes les preuves ne se valent pas. Une preuve est d'autant plus fiable
qu'elle est obtenue directement par l'auditeur, qu'elle est produite par
un système plutôt que rédigée par une personne, et qu'elle vient d'une
source indépendante de l'audité.

\begin{tableaucours}{cL{5.4cm}L{5.8cm}}{%
    \tableauHeader{Rang} &
    \tableauHeader{Type de preuve} &
    \tableauHeader{Exemple dans la mission}
}
1 & Obtenue directement par l'auditeur (test, ré-exécution, observation)
& Essai de connexion sans second facteur, restauration observée \\
2 & Produite par le système, extraite sous les yeux de l'auditeur &
Extraction des comptes, du collecteur de journaux \\
3 & Document de tiers indépendant & Rapport du test d'intrusion A3 \\
4 & Document produit par l'audité & Rapport de la RSSI, captures d'écran
fournies \\
5 & Déclaration orale & « Toutes les agences sont raccordées » \\
\end{tableaucours}
\vspace{-0.8em}
\begin{center}{\small\itshape Tableau~6.4. Une hiérarchie de la fiabilité
des preuves}\end{center}

\begin{boxinfo}
\textbf{Suffisante, pertinente, fiable.} Une preuve doit être
\emph{pertinente} (elle porte sur le critère), \emph{fiable} (rang élevé
du Tableau~6.4) et \emph{suffisante} (en quantité, pour la conclusion
visée). Une déclaration orale peut orienter un test~; elle ne suffit
jamais à coter \OUI{}.
\end{boxinfo}

\subsection{Échantillonner}

On ne vérifie pas toujours tout. Pour un \emph{contrôle automatisé} (le
second facteur, géré par un serveur central), on extrait l'ensemble de la
population~: 1\,240 comptes, sans échantillon. Pour ce qui se fait agence
par agence, l'équipe retient un échantillon de huit agences sur
trente-huit~:
\begin{itemize}
    \item Ngaoundéré, point d'entrée de l'attaque de 2026~;
    \item deux autres agences touchées en 2026~;
    \item les trois agences qui traitent le plus d'opérations~;
    \item deux agences tirées au hasard parmi les autres, dont Kribi.
\end{itemize}

\begin{boxdef}
\termecle{Échantillon raisonné}\index{échantillonnage}~: échantillon
choisi par l'auditeur selon les risques. Il permet de trouver des défauts
là où ils sont probables~; il ne permet pas d'extrapoler un taux à la
population. Seul un \termecle{échantillon aléatoire}, de taille suffisante,
autorise une estimation statistique.
\end{boxdef}

La distinction n'est pas un détail~: si trois agences de l'échantillon sur
huit présentent un défaut, on ne peut pas en déduire que «~37,5~\% des
agences~» le présentent. On peut seulement dire que le défaut existe, et
demander une vérification exhaustive.

% ============================================================================
\section{Coter les Tests}
\label{sec:c06-coter}
% ============================================================================

\subsection{Les résultats de terrain}

\begin{tableaucours}{cL{6.8cm}cL{2.8cm}}{%
    \tableauHeader{Réf.} &
    \tableauHeader{Résultat} &
    \tableauHeader{Cote} &
    \tableauHeader{Rang de la preuve}
}
T1 & 1\,240 comptes sur 1\,240 avec second facteur~; connexion refusée sans
lui dans les huit agences & \OUI{} & 1 et 2 \\
T2 & Tous les essais d'accès entre agences refusés & \OUI{} & 1 \\
T3 & Trois restaurations réussies, empreintes identiques & \OUI{} & 1 \\
T4 & Garoua sert en mode dégradé en 40 minutes~; Kribi ne peut pas servir
pendant plus de deux heures (procédure introuvable, codes inconnus) &
\NON{} & 1 \\
T5 & Listes de 30 agences, signatures non nominatives~; rapprochement
avec les effectifs impossible & \NEK{} & 4 \\
T6 & 31 agences sur 38 envoient des journaux au collecteur & \NON{} & 2 \\
T7 & Chef d'agence encore administrateur local dans 3 agences sur 8 &
\NON{} & 2 \\
T8 & Le prestataire des GAB refuse l'accès aux configurations, en
invoquant son contrat & \NEK{} & -- \\
\end{tableaucours}
\vspace{-0.8em}
\begin{center}{\small\itshape Tableau~6.5. Résultats des tests (mars 2028)}\end{center}

\begin{boxlecture}{La mission sur le réseau d'agences}
\textbf{Constatation.} Sur huit tests~: 3~\OUI{}, 3~\NON{}, 2~\NEK{}.
\[
\mu = 0{,}375~;\quad U = 0{,}25~;\quad
H = -(2 \times 0{,}375\log_2 0{,}375 + 0{,}25\log_2 0{,}25)
  \approx 0{,}531 + 0{,}531 + 0{,}500 \approx 1{,}56~\text{bit}.
\]

\textbf{Analyse} (confiance élevée sur les \OUI{} et les \NON{}). Les trois
mesures techniques centrales du plan (T1 à T3) fonctionnent~: le résiduel
de R-AG-01 sur l'entrée et la propagation est confirmé. Mais le rapport de
la RSSI déclarait «~tout en place~», et trois mesures ne le sont pas. Ce
n'est pas un défaut d'honnêteté~: la RSSI a reçu des déclarations des
agences et de la DSI, et les a transmises. C'est exactement la raison
d'être de l'audit~: remplacer des déclarations de rang~4 ou~5 par des
preuves de rang~1 ou~2.

\medskip
\hfill\lecture{0{,}375}{0{,}25}{1{,}56}
\end{boxlecture}

\subsection{Le cas du NEK}

Les deux \NEK{} ne sont pas de même nature.

\begin{itemize}
    \item \textbf{T5} est un \NEK{} de \emph{preuve insuffisante}~: les
          listes existent, mais ne permettent pas de conclure. L'auditeur
          recommande une preuve meilleure (émargement nominatif, quiz
          enregistré) et pourra conclure au prochain passage.
    \item \textbf{T8} est un \NEK{} de \emph{limitation de l'étendue}~: on
          a empêché l'auditeur de voir. C'est en soi un constat, grave, que
          le rapport porte au comité~: quatorze agences ont des GAB sous un
          système sans correctifs (Chapitre~\ref{chap:iso27005}), et la
          banque ne peut pas vérifier ce que fait leur prestataire.
\end{itemize}

\begin{boxalert}
\textbf{Violation typique.} Retirer du rapport un test que l'on n'a pas pu
réaliser, pour ne présenter que ce qui a été vérifié.

\smallskip
\textbf{Règle~:} une limitation de l'étendue figure dans le rapport, avec
sa cause et ce qu'elle empêche de conclure. Le principe de présentation
impartiale l'exige~: le commanditaire doit savoir ce que l'audit n'a pas
pu voir.
\end{boxalert}

\subsection{Le papier de travail}

\begin{boxprocedure}
\textbf{Le papier de travail d'un test.} Chaque test du programme laisse
un papier de travail qui permet à un autre auditeur de refaire le
raisonnement~:
\begin{itemize}
    \item la référence du test, le critère, la population et l'échantillon,
          avec la méthode de sélection~;
    \item la technique employée et la date~;
    \item les preuves obtenues, leur rang de fiabilité, et l'endroit où
          elles sont conservées (extraction, empreinte, procès-verbal)~;
    \item le résultat et la cote (\OUI{}, \NON{} ou \NEK{}), avec, pour un
          \NEK{}, sa nature (preuve insuffisante ou limitation)~;
    \item l'auteur et le réviseur.
\end{itemize}
\end{boxprocedure}

% ============================================================================
\section{Le Contradictoire et le Rapport}
\label{sec:c06-contradictoire}
% ============================================================================

\subsection{La réunion de clôture}

Avant le rapport, l'équipe présente ses constats aux audités~: la RSSI, la
DSI, la direction du réseau d'agences. C'est le
\termecle{contradictoire}\index{contradictoire}~: l'audité peut contester
un fait, apporter une preuve que l'auditeur n'avait pas, ou expliquer une
cause. Un constat qui a survécu au contradictoire est solide~; un constat
qui le contourne est fragile.

\begin{boxscenario}{Opération SANAGA. La DSI conteste T6}
\textbf{Faits.} À la réunion de clôture, le DSI affirme que les sept
agences manquantes «~ont été raccordées la semaine dernière~».

\medskip
\textbf{Réponse de l'équipe.} Mme~EKAMBI demande une extraction datée du
collecteur, faite devant un auditeur. L'extraction, obtenue le lendemain,
montre 38 agences sur 38.

\medskip
\textbf{Analyse} (confiance élevée). Le constat n'est pas supprimé~: pendant
la période auditée, sept agences n'étaient pas journalisées, et une
attaque y serait passée inaperçue. Le rapport le maintient, et ajoute
«~corrigé en cours de mission, extraction du 22~mars 2028~». La correction
est une bonne nouvelle~; elle ne réécrit pas le passé.
\end{boxscenario}

\subsection{Le rapport et le suivi}

Le rapport présente les objectifs, le périmètre, les critères, les
limitations, puis les constats et les conclusions. Chaque constat donne le
critère, le fait, la cote et la preuve~; le Chapitre~\ref{chap:conclure}
montre comment le développer en observation complète (cause, conséquence,
recommandation) et comment tirer une conclusion d'ensemble des cotes.
L'audité répond par un plan d'actions daté~; l'audit interne vérifie, à
l'échéance, que les actions sont faites, par des preuves du même rang que
celles de la mission.

\begin{boxobjet}{R-AG-01 · Après l'audit de mars 2028}
\textbf{Mesures confirmées par l'audit~:} second facteur (T1),
cloisonnement (T2), sauvegardes restaurées (T3).

\textbf{Mesures non confirmées~:} mode dégradé (T4, \NON{} à Kribi),
sensibilisation (T5, \NEK{}), journalisation (T6, corrigée en cours de
mission), administration locale (T7, \NON{} dans 3 agences sur 8).

\textbf{Résiduel~:} inchangé à $2 \times 3 = 6$ selon la convention de
réduction, car les étapes entrée, propagation et récupération restent
couvertes par des mesures prouvées. Mais la récupération repose sur les
sauvegardes seules~: dans une agence comme Kribi, les guichets resteraient
fermés le temps de la restauration.

\textbf{Nouveau risque à apprécier~:} compromission d'un GAB (T8), base
\NEK{}~; présentation au comité des risques en raison de la limitation de
l'étendue.
\end{boxobjet}

% ============================================================================
\section{Le Cadre Camerounais}
\label{sec:c06-cadre}
% ============================================================================

\begin{boxjuris}
\textbf{\loicam{n°~2010/012}, article~4, point~9}~\cite{loi2010012}~:

\medskip
\textit{« Audit de sécurité~: examen méthodique des composantes et des
acteurs de la sécurité, de la politique, des mesures, des solutions, des
procédures et des moyens mis en oeuvre par une organisation, pour
sécuriser son environnement, effectuer des contrôles de conformité, des
contrôles d'évaluation de l'adéquation des moyens (organisationnels,
techniques, humains, financiers) investis au regard des risques encourus,
d'optimisation, de rationalité et de performance. »}

\medskip
\textbf{Article~13, alinéa~1}~:

\medskip
\textit{« Sont soumis à un audit de sécurité obligatoire, les réseaux de
communications électroniques et les systèmes d'information, des
opérateurs, des autorités de certification et des fournisseurs de services
de communications électroniques. »}

\medskip
\textbf{Article~32, alinéas~1 et~2}~:

\medskip
\textit{« (1) Les réseaux de communications électroniques et les systèmes
d'information sont soumis à un régime d'audit de sécurité obligatoire et
périodique de leurs systèmes de sécurité par l'Agence. (2) L'audit de
sécurité et les mesures d'impact de gravité sont effectués chaque année ou
lorsque les circonstances l'exigent. »}

\medskip
\textbf{Portée.} La définition légale de l'article~4 couvre à la fois la
conformité et l'adéquation des moyens «~au regard des risques encourus~»~:
elle correspond à l'approche fondée sur le risque de l'ISO~19011. Les
articles~13 et~32 instituent un audit réglementaire annuel, conduit sous
l'autorité de l'Agence nationale des technologies de l'information et de la
communication. Leur champ exact (quels systèmes, quelles organisations)
et leur application à une agence publique comme l'ANPP sont étudiés au
Chapitre~\ref{chap:droit}.
\end{boxjuris}

% ============================================================================
\section{Ce que l'Audit demande à chaque Métier}
\label{sec:c06-metiers}
% ============================================================================

\begin{boxmetier}{Auditeur en cybersécurité}
\textbf{Préférer la preuve que l'on obtient soi-même}
\begin{enumerate}
    \item Obtenir un mandat écrit qui garantit l'accès, y compris chez les
          prestataires.
    \item Écrire, avant le terrain, la preuve attendue de chaque test.
    \item Viser les preuves de rang~1 et~2~; ne jamais coter \OUI{} sur une
          déclaration.
    \item Distinguer, pour chaque \NEK{}, preuve insuffisante et limitation
          de l'étendue, et porter la seconde au commanditaire.
    \item Maintenir un constat corrigé en cours de mission, en signalant la
          correction.
\end{enumerate}
\end{boxmetier}

\begin{boxmetier}{RSSI}
\textbf{Préparer l'audit, pas le subir}

La RSSI aurait évité l'écart entre son rapport et l'audit en demandant,
pour chaque action, la preuve prévue au plan (Chapitre~\ref{chap:iso27005})
plutôt qu'une déclaration. Un RSSI dont les rapports reposent sur des
preuves de rang~2 voit l'audit confirmer ses chiffres.
\end{boxmetier}

\begin{boxmetier}{Responsable GRC}
\textbf{Faire entrer l'audit dans les contrats}

La limitation de T8 vient d'un contrat sans droit d'audit, comme celui de
DIGICORE. Le responsable GRC fait inscrire, dans chaque contrat de
prestataire qui touche au système d'information, un droit d'accès de
l'audit interne et des auditeurs mandatés.
\end{boxmetier}

\separateur

\begin{boxresume}
\textbf{Ce qu'il faut retenir de ce chapitre~:}
\begin{itemize}
    \item \textbf{L'audit} obtient et évalue des preuves au regard de
          critères~; le constat peut être conforme, non conforme, ou
          \NEK{}.
    \item \textbf{Sept principes} (ISO~19011)~: intégrité, présentation
          impartiale, conscience professionnelle, confidentialité,
          indépendance, approche fondée sur la preuve, approche fondée sur
          le risque.
    \item \textbf{Quatre familles}~: interne, contractuel, certification,
          réglementaire.
    \item \textbf{Le cycle}~: mandat, préparation, terrain, contradictoire,
          rapport, suivi.
    \item \textbf{La preuve} vaut par sa pertinence, sa fiabilité et sa
          suffisance~; une déclaration n'est jamais suffisante.
    \item \textbf{Un échantillon raisonné} trouve des défauts~; il ne
          permet pas d'extrapoler un taux.
    \item \textbf{Une limitation de l'étendue} est un constat, porté au
          commanditaire.
\end{itemize}
\end{boxresume}

\begin{boxinfo}
\textbf{Points de contrôle.} Vous avez compris ce chapitre si vous savez~:
\begin{itemize}
    \item classer une preuve dans la hiérarchie de fiabilité~;
    \item rédiger une ligne de programme de travail avec sa preuve
          attendue~;
    \item dire ce qu'un échantillon raisonné permet de conclure~;
    \item distinguer les deux sortes de \NEK{} et traiter chacune.
\end{itemize}
\end{boxinfo}

% ============================================================================
\section*{Exercices}
% ============================================================================

\begin{boxexo}[title={\ding{45}\quad Exercice~6.1. Techniques et preuves \niveaubase}]
Pour chacun des éléments suivants, indiquez la technique (examiner,
interroger, tester, observer, ré-exécuter) et le rang de fiabilité de la
preuve (Tableau~6.4).
\begin{enumerate}[(a)]
    \item Le DSI explique que les correctifs sont appliqués chaque mois.
    \item L'auditeur lance lui-même la restauration d'un serveur.
    \item La RSSI remet une capture d'écran de la console de sauvegarde.
    \item L'auditeur assiste à l'ouverture d'une agence en mode dégradé.
    \item L'auditeur extrait la liste des comptes d'administration, devant
          lui, depuis l'annuaire.
    \item L'auditeur recalcule le délai d'application des correctifs à
          partir du journal de déploiement.
\end{enumerate}
\end{boxexo}

\begin{boxexo}[title={\ding{45}\quad Exercice~6.2. Une ligne de programme \niveaumoyen}]
Le comité demande d'ajouter à la mission un test sur les comptes des agents
partis~: aucun compte d'un agent ayant quitté la banque ne doit rester
actif plus de cinq jours après son départ.
\begin{enumerate}[(a)]
    \item Rédigez la ligne du programme de travail~: critère, test,
          population, échantillon éventuel, preuve attendue.
    \item La DRH fournit une liste de 46 départs en 2027~; l'extraction de
          l'annuaire montre 4 comptes de partants actifs plus de cinq jours.
          Cotez le test et justifiez.
    \item La DRH refuse de communiquer les départs du dernier trimestre,
          «~pour des raisons de confidentialité~». Quelle est la nature de
          ce \NEK{}, et que fait l'auditeur~?
\end{enumerate}
\end{boxexo}

\begin{boxexo}[title={\ding{45}\quad Exercice~6.3. Ce que dit l'échantillon \niveaumoyen}]
\begin{enumerate}[(a)]
    \item Justifiez chacun des quatre critères de sélection de l'échantillon
          de huit agences.
    \item Un administrateur lit le résultat de T7 et conclut que «~14
          agences environ~» gardent l'administration locale. Que
          répondez-vous~?
    \item Si les huit agences avaient été conformes pour T7, l'auditeur
          aurait-il pu coter \OUI{} pour les 38 agences~? Proposez la
          meilleure preuve possible pour ce test.
\end{enumerate}
\end{boxexo}

\begin{boxexo}[title={\ding{45}\quad Exercice~6.4. L'audit contractuel de PHARMACAM \niveauexpert}]
PHARMACAM exerce, en 2028, son droit d'audit sur LiZbiZ, pour vérifier le
plan de traitement du scénario SO1 (Chapitres~\ref{chap:ebios}
et~\ref{chap:fair}).
\begin{enumerate}[(a)]
    \item De quelle famille d'audit s'agit-il~? Qui peut le conduire~?
          AUDIX CONSEIL peut-il s'en charger~? Justifiez.
    \item Rédigez la lettre de mission~: objectif, périmètre, critères,
          accès.
    \item Rédigez un programme de travail de cinq tests, avec la preuve
          attendue et son rang.
    \item L'éditeur de l'application refuse de montrer la configuration de
          son accès de maintenance. Comment le rapport le traite-t-il, et
          quelle conséquence PHARMACAM peut-elle en tirer~?
\end{enumerate}
\end{boxexo}

\bigskip

\begin{boxinfo}
\textbf{Pour aller plus loin.}
\begin{itemize}
    \item L'ISO~19011:2018~\cite{iso19011_2018}, sur les principes, le
          programme et la conduite des audits.
    \item L'ISO/IEC~27007:2020~\cite{iso27007_2020}, pour l'audit d'un SMSI.
    \item Le NIST SP~800-53A~\cite{nist2022sp80053a}, sur les méthodes
          d'évaluation des mesures de sécurité.
    \item Les normes mondiales d'audit interne de l'IIA~\cite{iia2024gias}.
\end{itemize}
\end{boxinfo}

\begin{boxinfo}
\textbf{Transition.} La mission a produit huit cotes~: 3~\OUI{}, 3~\NON{},
2~\NEK{}. Que conclure~? Le réseau d'agences est-il « sécurisé »~? Pour
répondre, il faut décomposer l'objet audité, coter chaque assertion selon
plusieurs critères, et appliquer une règle de décision qui tienne compte de
l'incertitude. C'est l'objet du Chapitre~\ref{chap:conclure}.
\end{boxinfo}
